\section{Introduction}

Where does a large language model store its facts? In this paper, we report evidence that factual associations in GPT correspond to a localized computation that can be directly edited.


Large language models can predict factual statements about the world~\citep{petroni-etal-2019-language,jiang-etal-2020-know,roberts-etal-2020-much}. For example, given the prefix ``\emph{The Space Needle is located in the city of},'' GPT will reliably predict the true answer: ``\emph{Seattle}'' (\reffig{infoflow}a).  Factual knowledge has been observed to emerge in both autoregressive GPT models~\citep{gpt2,gpt3} and masked BERT models~\citep{devlin-etal-2019-bert}.

In this paper, we investigate how such factual associations are stored within GPT-like autoregressive transformer models. Although many of the largest neural networks in use today are autoregressive, the way that they store knowledge remains under-explored. Some research has been done for masked models~\citep{petroni-etal-2019-language,jiang-etal-2020-know,10.1162/tacl_a_00410,geva-etal-2021-transformer,dai-kn,decao-ke}, but GPT has architectural differences such as unidirectional attention and generation capabilities that provide an opportunity for new insights.

We use two approaches. First, we trace the causal effects of hidden state \underline{activations} within GPT using causal mediation analysis~\citep{pearl2001direct,vig2020investigating} to identify the specific modules that mediate recall of a fact about a subject (\reffig{infoflow}).  Our analysis reveals that feedforward MLPs at a range of middle layers are decisive when processing the last token of the subject name (Figures~\ref{fig:infoflow}b,\ref{fig:avg-1000-trace}b,\ref{fig:trace-mlp-disabled}).

Second, we test this finding in model \underline{weights} by introducing a Rank-One Model Editing method (ROME) to alter the parameters that determine a feedfoward layer's behavior at the decisive token.  Despite the simplicity of the intervention, we find that ROME is similarly effective to other model-editing approaches on a standard zero-shot relation extraction benchmark (Section~\ref{subsec:zsre}).

To evaluate ROME's impact on more difficult cases, we introduce a dataset of counterfactual assertions (Section~\ref{subsec:counterfact-description}) that would not have been observed in pretraining. Our evaluations (Section~\ref{subsec:romeeval-mlp}) confirm that midlayer MLP modules can store factual associations that generalize beyond specific surface forms, while remaining specific to the subject. Compared to previous fine-tuning~\citep{zhu-ft}, interpretability-based~\citep{dai-kn}, and meta-learning~\citep{mend,decao-ke} methods, ROME achieves good generalization and specificity simultaneously, whereas previous approaches sacrifice one or the other.
